% !TeX root = ../sections/02_exercises.tex

\begin{exercise}[情報数学 A 中間試験 1]
	周期 \(2\pi\) をもち，区間 \([-\pi,\pi)\) 上で
	\[
	f(x)=x
	\]
	で与えられる周期関数 \(f(x)\) のフーリエ級数を求めよ．
	
	\medskip
	
	※解答は \(5x\) の項まで書き下せ．
	例えば，周期 \(2\pi\) をもち，区間 \([-\pi,\pi)\) 上で
	\[
	g(x)=|x|
	\]
	で与えられる周期関数 \(g(x)\) のフーリエ係数は
	\[
	a_0=\pi,\qquad
	a_n=
	\begin{cases}
		-\dfrac{4}{\pi n^2} & (n \text{ が奇数}),\\[6pt]
		0 & (n \text{ が偶数}),
	\end{cases}
	\qquad
	b_n=0
	\]
	である．この場合，\(g(x)\) のフーリエ級数は次のように書き表す：
	\[
	g(x)\sim
	\frac{\pi}{2}
	-\frac{4}{\pi}
	\left(
	\cos x+\frac{1}{3^2}\cos 3x+\frac{1}{5^2}\cos 5x+\cdots
	\right).
	\]
\end{exercise}
\begin{background}
	この問題では，周期 \(2\pi\) の関数
	\[
	f(x)=x
	\]
	のフーリエ級数を求める。
	
	周期 \(2\pi\) の関数のフーリエ級数は，
	\[
	f(x)\sim \frac{a_0}{2}
	+\sum_{n=1}^{\infty}
	\left(a_n\cos nx+b_n\sin nx\right)
	\]
	の形で表される。
	
	ここで，フーリエ係数は
	\[
	a_0=\frac{1}{\pi}\int_{-\pi}^{\pi}f(x)\,dx,
	\]
	\[
	a_n=\frac{1}{\pi}\int_{-\pi}^{\pi}f(x)\cos nx\,dx,
	\]
	\[
	b_n=\frac{1}{\pi}\int_{-\pi}^{\pi}f(x)\sin nx\,dx
	\]
	で定義される。
	
	この問題では，\(f(x)=x\) が奇関数であることが重要である。
	つまり，
	\[
	f(-x)=-f(x)
	\]
	が成り立つ。
	
	また，
	\[
	\cos nx
	\]
	は偶関数であり，
	\[
	\sin nx
	\]
	は奇関数である。
	
	したがって，
	\[
	x\cos nx
	\]
	は奇関数であり，
	\[
	x\sin nx
	\]
	は偶関数である。
	
	奇関数を対称区間 \([-\pi,\pi]\) で積分すると \(0\) になるため，
	\[
	a_0=0,\qquad a_n=0
	\]
	である。
	
	よって，この問題では正弦係数 \(b_n\) だけを求めればよい。
	これは，奇関数のフーリエ級数が正弦級数になるという典型例である。
\end{background}
\begin{strategy}
	この問題を見たら，まず
	\[
	f(x)=x
	\]
	が奇関数であることに注目する。
	
	奇関数なので，フーリエ級数には余弦項が現れない。
	したがって，
	\[
	a_0=0,\qquad a_n=0
	\]
	と分かる。
	
	残るのは
	\[
	b_n=\frac{1}{\pi}\int_{-\pi}^{\pi}x\sin nx\,dx
	\]
	の計算である。
	
	ここで \(x\sin nx\) は偶関数なので，
	\[
	b_n=\frac{2}{\pi}\int_0^\pi x\sin nx\,dx
	\]
	とできる。
	
	この積分は部分積分で計算する。
	\[
	u=x,\qquad dv=\sin nx\,dx
	\]
	とおくと，
	\[
	du=dx,\qquad v=-\frac{\cos nx}{n}
	\]
	である。
	
	したがって，
	\[
	\int_0^\pi x\sin nx\,dx
	=
	\left[-\frac{x\cos nx}{n}\right]_0^\pi
	+
	\frac{1}{n}\int_0^\pi \cos nx\,dx
	\]
	となる。
	
	第2項は
	\[
	\int_0^\pi \cos nx\,dx
	=
	\left[\frac{\sin nx}{n}\right]_0^\pi
	=0
	\]
	である。
	
	よって，
	\[
	\int_0^\pi x\sin nx\,dx
	=
	-\frac{\pi\cos n\pi}{n}
	=
	-\frac{\pi(-1)^n}{n}
	\]
	である。
	
	したがって，
	\[
	b_n
	=
	\frac{2}{\pi}
	\left(
	-\frac{\pi(-1)^n}{n}
	\right)
	=
	\frac{2(-1)^{n+1}}{n}
	\]
	となる。
	
	以上より，
	\[
	x\sim
	2\sum_{n=1}^{\infty}
	\frac{(-1)^{n+1}}{n}\sin nx
	\]
	である。
	
	問題文では \(5x\) の項まで書くように求められているので，
	最後は
	\[
	x\sim
	2\left(
	\sin x-\frac{1}{2}\sin 2x
	+\frac{1}{3}\sin 3x
	-\frac{1}{4}\sin 4x
	+\frac{1}{5}\sin 5x
	+\cdots
	\right)
	\]
	のように書き下す。
	
	この問題の発想は，
	\[
	\text{奇関数}
	\Rightarrow
	\text{正弦項だけ}
	\Rightarrow
	b_n\text{ を部分積分で求める}
	\]
	という流れである。
\end{strategy}